\section{Conclusão}
\label{sec:conclusao}

Este trabalho em andamento estruturou uma investigação aplicada para prevenir a
liberação de combustível incompatível antes do abastecimento. A pergunta,
hipóteses e objetivos foram convertidos em um método que combina DSRM,
mapeamento da literatura e estudo de caso, além de requisitos, arquitetura,
estados e critérios de aceitação rastreáveis.

A análise documental parcial confirmou que os campos de combustível pertencem à
consulta detalhada de veículo e que o acesso envolve autorização SENATRAN,
Serpro e GCC. Também corrigiu a premissa de preço: R\$~0{,}80 não é o valor da
primeira faixa da consulta detalhada na Portaria n\textsuperscript{o} 417/2026. As restrições de
segurança levaram a manter o ESP32 e os atuadores comuns exclusivamente em
bancada seca e a considerar controladores certificados como caminho para uma
eventual aplicação industrial.

O primeiro incremento de software materializou backend, PWA e firmware
experimental em modo sintético. A verificação local aprovou 82 testes do
backend, com 93\% de cobertura de instruções, 12 testes da interface, a
compilação de produção da PWA, o teste em computador da lógica de
intertravamento do firmware e o build PlatformIO para o alvo \path{esp32dev}.
Esses resultados não incluem integração MQTT, gravação/execução em placa ou
ensaio físico.

Ainda não há base para concluir que o sistema reduz incidentes ou é
economicamente viável. Dados da empresa, entrevistas, matriz integral de
software e ensaios de hardware permanecem pendentes; logo, H1 e H2 não foram
testadas conclusivamente. A próxima etapa é executar o protocolo de revisão,
obter as autorizações, integrar os componentes simulados, validar a bancada e
produzir os resultados experimentais sem ultrapassar os limites éticos,
regulatórios e físicos definidos.
